\ifdefined\gkver\else\def\gkver{student}\fi
\documentclass[\gkver]{gaokaozhenti}
\begin{document}

\chapter{2016年4月浙江}

\begin{enumerate}
\item
%% number: 1
%% paperName: 2016年4月浙江选考第 1 题 3 分
%% typeId: 1
%% type: 单选题
%% chapter: 直线运动
%% point: 路程与位移
%% method: 无
%% score: 3
%% degree: 950
%% duplicateId: 0
%% body:
某同学绕操场一周跑了 $400\Um$，用时 $65\Us$。这两个物理量分别是\xzanswer{C}

\fourchoices[answer=C]
{路程、时刻}
{位移、时刻}
{路程、时间}
{位移、时间}

%% answer: C
%% memo:
\memoanswer{
	绕 $400\Um$ 操场跑一圈，首末位置重合，则位移的大小为 $0$，路程等于 $400\Um$。$65\Us$ 指时间长度，是指时间，故 C 正确。
}

\item
%% number: 2
%% paperName: 2016年4月浙江选考第 2 题 3 分
%% typeId: 1
%% type: 单选题
%% chapter: 曲线运动
%% point: 平抛运动
%% method: 无
%% score: 3
%% degree: 950
%% duplicateId: 0
%% body:
如图所示是某人在投飞镖，飞镖在飞行途中受到的力有\xzanswer{B}

\onepicture[h=3.5cm]{figs/02.png}

\fourchoices[answer=B]
{推力}
{重力、空气阻力}
{重力、推力}
{重力、推力、空气阻力}

%% answer: B
%% memo:
\memoanswer{
	飞镖在飞行途中受到的力有重力、空气阻力，没有向前的推力，飞镖离开手，不再受到力的作用，故 B 正确。
}

\item
%% number: 3
%% paperName: 2016年4月浙江选考第 3 题 3 分
%% typeId: 1
%% type: 单选题
%% chapter: 直线运动
%% point: 参考系
%% method: 无
%% score: 3
%% degree: 950
%% duplicateId: 0
%% body:
下列说法正确的是\xzanswer{B}

\fourchoices[answer=B]
{物体在完全失重的状态下没有惯性}
{运动是绝对的，但运动的描述是相对的}
{电流强度有大小和方向，所以是矢量}
{研究月球绕地球运行轨迹时不能把月球看成质点}

%% answer: B
%% memo:
\memoanswer{
	物体在完全失重的状态下因为质量不变，则惯性不变，故 A 错误；运动是绝对的，但运动的描述是相对的，故 B 正确；电流有大小和方向，但是电流的合成不满足平行四边形定则，所以是标量，故 C 错误；研究月球绕地球运行轨迹时，月球的大小和形状可忽略不计，则可把月球看成质点，故 D 错误。
}

\item
%% number: 4
%% paperName: 2016年4月浙江选考第 4 题 3 分
%% typeId: 1
%% type: 单选题
%% chapter: 其他
%% point: 物理学史
%% method: 无
%% score: 3
%% degree: 950
%% duplicateId: 0
%% body:
物理学中的自由落体规律、万有引力定律、静止点电荷之间的相互作用规律和电流磁效应分别由不同的物理学家探究发现，他们依次是\xzanswer{A}

\fourchoices[answer=A]
{伽利略、牛顿、库仑和奥斯特}
{牛顿、安培、洛伦兹和奥斯特}
{伽利略、卡文迪许、库仑和安培}
{开普勒、伽利略、库仑和洛伦兹}

%% answer: A
%% memo:
\memoanswer{
	物理学中的自由落体规律、万有引力定律、静止点电荷之间的相互作用规律和电流磁效应分别由不同的物理学家探究发现，他们依次是：伽利略、牛顿、库仑和奥斯特，故 A 正确，BCD 错误。
}

\item
%% number: 5
%% paperName: 2016年4月浙江选考第 5 题 3 分
%% typeId: 1
%% type: 单选题
%% chapter: 曲线运动
%% point: 圆周运动的应用
%% method: 无
%% score: 3
%% degree: 850
%% duplicateId: 0
%% body:
如图为某中国运动员在短道速滑比赛中勇夺金牌的精彩瞬间。假定此时他正沿圆弧形弯道匀速率滑行，则他\xzanswer{D}

\onepicture[h=3.5cm]{figs/05.jpg}

\fourchoices[answer=D]
{所受的合力为零，做匀速运动}
{所受的合力恒定，做匀加速运动}
{所受的合力恒定，做变加速运动}
{所受的合力变化，做变加速运动}

%% answer: D
%% memo:
\memoanswer{
	匀速圆周运动过程中，线速度大小不变，方向改变，向心加速度大小不变，方向始终指向圆心，向心力大小不变，方向始终指向圆心。所以物体做变加速运动，故 D 正确，ABC 错误。
}

\item
%% number: 6
%% paperName: 2016年4月浙江选考第 6 题 3 分
%% typeId: 1
%% type: 单选题
%% chapter: 直线运动
%% point: 自由落体运动
%% method: 无
%% score: 3
%% degree: 950
%% duplicateId: 0
%% body:
宇航员在月球上离月球表面高 $10\Um$ 处由静止释放一片羽毛，羽毛落到月球表面上的时间大约是\xzanswer{C}

\fourchoices[answer=C]
{$1.0\Us$}
{$1.4\Us$}
{$3.5\Us$}
{$12\Us$}

%% answer: C
%% memo:
\memoanswer{
	月球表面的加速度约为地球表面的加速度的 $\frac{1}{6}$，由 $h=\frac{1}{2}g_{\text{月}}t^{2}$ 解得 $t=\sqrt{\frac{2h}{g_{\text{月}}}}=\sqrt{\frac{2\times10\times6}{10}}=\sqrt{12}\Us\approx3.5\Us$，故 C 正确，ABD 错误。
}

\item
%% number: 7
%% paperName: 2016年4月浙江选考第 7 题 3 分
%% typeId: 1
%% type: 单选题
%% chapter: 电路
%% point: 电容
%% method: 无
%% score: 3
%% degree: 950
%% duplicateId: 0
%% body:
关于电容器，下列说法正确的是\xzanswer{D}

\fourchoices[answer=D]
{在充电过程中电流恒定}
{在放电过程中电容减小}
{能储存电荷，但不能储存电能}
{两个彼此绝缘又靠近的导体可视为电容器}

%% answer: D
%% memo:
\memoanswer{
	电容器在充电过程中电流是变化的，故 A 错误；在放电过程中，电容的大小是由电容器的结构决定的，不会发生变化，故 B 错误；电容器可以储存电荷量也就是可以储存电能，故 C 错误；两个彼此绝缘相互靠近的导体可以构成一个电容器，故 D 正确。
}

\item
%% number: 8
%% paperName: 2016年4月浙江选考第 8 题 3 分
%% typeId: 1
%% type: 单选题
%% chapter: 电场
%% point: 带电粒子的平衡
%% method: 无
%% score: 3
%% degree: 850
%% duplicateId: 0
%% body:
密立根油滴实验原理如图所示。两块水平放置的金属板分别与电源的正负极相接，板间电压为 $U$，形成竖直向下场强为 $E$ 的匀强电场。用喷雾器从上板中间的小孔喷入大小、质量和电荷量各不相同的油滴。通过显微镜可找到悬浮不动的油滴，若此悬浮油滴的质量为 $m$，则下列说法正确的是\xzanswer{C}

\onepicture[w=6.0cm]{figs/08.png}

\fourchoices[answer=C]
{悬浮油滴带正电}
{悬浮油滴的电荷量为 $\frac{mg}{U}$}
{增大场强，悬浮油滴将向上运动}
{油滴的电荷量不一定是电子电量的整数倍}

%% answer: C
%% memo:
\memoanswer{
	带电荷量为 $q$ 的油滴静止不动，则油滴受到向上的电场力，而题图中板间场强方向竖直向下，则油滴带负电，故 A 错误；由平衡条件得 $mg=q\frac{U}{d}$，解得 $q=\frac{mgd}{U}$，故 B 错误；由平衡条件得 $mg=qE$，当增大场强，电场力增大，则悬浮油滴将向上运动，故 C 正确；不同油滴的所带电荷量虽不相同，但都是某个最小电荷量（元电荷）的整数倍，故 D 错误。
}

\item
%% number: 9
%% paperName: 2016年4月浙江选考第 9 题 3 分
%% typeId: 1
%% type: 单选题
%% chapter: 磁场
%% point: 左手定则
%% method: 无
%% score: 3
%% degree: 650
%% duplicateId: 0
%% body:
法拉第电动机原理如图所示。条形磁铁竖直固定在圆形水银槽中心，N 极向上。一根金属杆斜插在水银中，杆的上端与固定在水银槽圆心正上方的铰链相连。电源负极与金属杆上端相连，与电源正极连接的导线插入水银中。从上往下看，金属杆\xzanswer{D}

\onepicture[h=4.0cm]{figs/09.png}

\fourchoices[answer=D]
{向左摆动}
{向右摆动}
{顺时针转动}
{逆时针转动}

%% answer: D
%% memo:
\memoanswer{
	电源、金属棒、导线和水银组成闭合电路，金属棒中有斜向上方的电流，由左手定则可知，导电金属杆所受安培力将会使其转动，从上往下看，金属杆将逆时针转动，故 D 正确。
}

\item
%% number: 10
%% paperName: 2016年4月浙江选考第 10 题 3 分
%% typeId: 1
%% type: 单选题
%% chapter: 曲线运动
%% point: 平抛运动
%% method: 无
%% score: 3
%% degree: 650
%% duplicateId: 0
%% body:
某卡车在公路上与路旁障碍物相撞。处理事故的警察在泥地中发现了一个小的金属物体，经判断，它是相撞瞬间车顶上一个松脱的零件被抛出而陷在泥里的。为了判断卡车是否超速，需要测量的量是\xzanswer{D}

\fourchoices[answer=D]
{车的长度，车的重量}
{车的高度，车的重量}
{车的长度，零件脱落点与陷落点的水平距离}
{车的高度，零件脱落点与陷落点的水平距离}

%% answer: D
%% memo:
\memoanswer{
	零件被抛出后做平抛运动，且平抛运动的初速度等于卡车原来的速度，由平抛规律得 $h=\frac{1}{2}gt^{2}$，$x=v_{0}t$，解得 $v_{0}=x\sqrt{\frac{g}{2h}}$，要测量卡车原来的速度，只要测出零件平抛运动的初速度，由上知，需要测量的量是车的高度 $h$，零件脱落点与陷落点的水平距离 $x$，故 D 正确。
}

\item
%% number: 11
%% paperName: 2016年4月浙江选考第 11 题 3 分
%% typeId: 1
%% type: 单选题
%% chapter: 曲线运动
%% point: 人造地球卫星
%% method: 无
%% score: 3
%% degree: 750
%% duplicateId: 0
%% body:
2015 年 12 月，我国暗物质粒子探测卫星“悟空”发射升空进入高为 $5.0\times10^{2}\Ukm$ 的预定轨道。“悟空”卫星和地球同步卫星的运动均可视为匀速圆周运动。已知地球半径 $R=6.4\times10^{3}\Ukm$。下列说法正确的是\xzanswer{C}

\onepicture[w=5.5cm]{figs/11.jpg}

\fourchoices[answer=C]
{“悟空”卫星的线速度比同步卫星的线速度小}
{“悟空”卫星的角速度比同步卫星的角速度小}
{“悟空”卫星的运行周期比同步卫星的运行周期小}
{“悟空”卫星的向心加速度比同步卫星的向心加速度小}

%% answer: C
%% memo:
\memoanswer{
	由 $G\frac{Mm}{r^{2}}=m\frac{v^{2}}{r}=m\omega^{2}r=m\left(\frac{2\pi}{T}\right)^{2}r=ma$ 得，$v=\sqrt{\frac{GM}{r}}$，$\omega=\sqrt{\frac{GM}{r^{3}}}$，$T=2\pi\sqrt{\frac{r^{3}}{GM}}$，$a=\frac{GM}{r^{2}}$；由 $v=\sqrt{\frac{GM}{r}}$ 可知半径小的线速度大，故 A 错误；由 $\omega=\sqrt{\frac{GM}{r^{3}}}$ 可知半径小的角速度大，故 B 错误；由 $T=2\pi\sqrt{\frac{r^{3}}{GM}}$ 可知半径小的周期小，故 C 正确；由 $a=\frac{GM}{r^{2}}$ 可知半径小的加速度大，故 D 错误。
}

\item
%% number: 12
%% paperName: 2016年4月浙江选考第 12 题 3 分
%% typeId: 1
%% type: 单选题
%% chapter: 电路
%% point: 电功电功率
%% method: 无
%% score: 3
%% degree: 650
%% duplicateId: 0
%% body:
图示中的路灯为太阳能路灯，每只路灯的光伏电池板有效采光面积约 $0.3\Umq$。晴天时电池板上每平方米每小时接收到的太阳辐射能约为 $3\times10^{6}\UJ$。如果每天等效日照时间约为 $6\Uh$，光电池一天产生的电能可供 $30\UW$ 的路灯工作 $8\Uh$。光电池的光电转换效率约为\xzanswer{C}

\onepicture[w=5.0cm]{figs/12.jpg}

\fourchoices[answer=C]
{$4.8\%$}
{$9.6\%$}
{$16\%$}
{$44\%$}

%% answer: C
%% memo:
\memoanswer{
	由题意知，每只路灯的有效采光面积 $S$，每平方米每小时接收太阳能 $E_{0}$，每天日照时间 $t$，则每只路灯每天接收到的太阳能为 $E_{\text{效}}=SE_{0}t$；光电池每天可以使 $P_{0}=30\UW$ 的路灯工作 $T_{0}=8\Uh$，则光电池光电转换率为 $\eta=\frac{P_{0}T_{0}}{SE_{0}t}\times100\%=16\%$，故 C 正确。
}

\item
%% number: 13
%% paperName: 2016年4月浙江选考第 13 题 3 分
%% typeId: 1
%% type: 单选题
%% chapter: 电场
%% point: 电场叠加
%% method: 无
%% score: 3
%% degree: 550
%% duplicateId: 0
%% body:
如图所示，真空中有两个点电荷 $Q_{1}=+9.0\times10^{-8}\UC$ 和 $Q_{2}=-1.0\times10^{-8}\UC$，分别固定在 $x$ 坐标轴上，其中 $Q_{1}$ 位于 $x=0$ 处，$Q_{2}$ 位于 $x=6\Ucm$ 处。在 $x$ 轴上\xzanswer{D}

\onepicture[w=6.0cm]{\includesvg{figs/13}}

\fourchoices[answer=D]
{场强为 $0$ 的点有两处}
{在 $x>6\Ucm$ 区域，电势沿 $x$ 轴正方向降低}
{质子从 $x=1\Ucm$ 运动到 $x=5\Ucm$ 处，电势能升高}
{在 $0<x<6\Ucm$ 和 $x>9\Ucm$ 的区域，场强沿 $x$ 轴正方向}

%% answer: D
%% memo:
\memoanswer{
	某点的电场强度是正电荷 $Q_{1}$ 和负电荷 $Q_{2}$ 在该处产生的电场的叠加，是合场强，根据点电荷的场强公式 $E=k\frac{Q}{r^{2}}$，所以要使电场强度为零，那么正电荷 $Q_{1}$ 和负电荷 $Q_{2}$ 在该处产生的场强必须大小相等、方向相反，不会在 $Q_{1}$ 的左边，因为 $Q_{1}$ 的电荷量大于 $Q_{2}$，也不会在 $Q_{1}Q_{2}$ 之间，因为它们电性相反，在中间的电场方向都向右，所以只能在 $Q_{2}$ 右边，即在 $x$ 坐标轴上电场强度为零的点只有一个，故 A 错误；由 A 的分析可知，在 $x>6\Ucm$ 区域，有一点使得电场强度为 $0$，则在 $x>6\Ucm$ 区域，沿 $x$ 轴正方向电势先升高后降低，故 B 错误；$0<x<6\Ucm$ 区域电场强度方向沿 $x$ 正方向，质子从 $x=1\Ucm$ 运动到 $x=5\Ucm$ 处，电势能减小，故 C 错误；在 $0<x<6\Ucm$ 的区域，场强沿 $x$ 轴正方向，当 $k\frac{Q_{1}}{x^{2}}=k\frac{Q_{2}}{(x-6)^{2}}$ 时，解得 $x=9\Ucm$，即 $x=9\Ucm$ 处场强为 $0$，则 $x>9\Ucm$ 处场强沿 $x$ 轴正方向，故 D 正确。
}

\item
%% number: 14
%% paperName: 2016年4月浙江选考第 14 题 2 分
%% typeId: 2
%% type: 多选题
%% chapter: 光学
%% point: 全反射
%% method: 无
%% score: 2
%% degree: 750
%% duplicateId: 0
%% body:
下列说法正确的是\xzanswer{BC}

\fourchoices[answer=BC]
{波源的频率越高，波速越大}
{温度升高，放射性元素的半衰期不变}
{氢原子吸收光子从低能级跃迁到高能级}
{光发生全反射时，临界角随入射角增大而变大}

%% answer: BC
%% memo:
\memoanswer{
	机械波的波速取决于介质，与频率无关，故 A 错误；半衰期是原子核本身的属性，与温度无关，故 B 正确；氢原子在吸收特定频率的光子后能从低能级跃迁到高能级，故 C 正确；由 $\sin C=\frac{1}{n}$ 可知，临界角只与材料的折射率有关，故 D 错误。
}

\item
%% number: 15
%% paperName: 2016年4月浙江选考第 15 题 2 分
%% typeId: 2
%% type: 多选题
%% chapter: 机械振动
%% point: 单摆
%% method: 无
%% score: 2
%% degree: 550
%% duplicateId: 0
%% body:
摆球质量相等的甲、乙两单摆悬挂点高度相同，其振动图象如图所示。选悬挂点所在水平面为重力势能的参考面，由图可知\xzanswer{AC}

\onepicture[w=6.0cm]{\includesvg{figs/15}}

\fourchoices[answer=AC]
{甲、乙两单摆的摆长之比是 $\frac{9}{4}$}
{$t_{a}$ 时刻甲、乙两单摆的摆角相等}
{$t_{b}$ 时刻甲、乙两单摆的势能差最大}
{$t_{c}$ 时刻甲、乙两单摆的速率相等}

%% answer: AC
%% memo:
\memoanswer{
	由振动图像得到甲、乙两个单摆的周期之比为 $T_{\text{甲}}:T_{\text{乙}}=2:3$，由单摆周期公式 $T=2\pi\sqrt{\frac{L}{g}}$ 得 $T^{2}\propto L$，故甲、乙两个单摆的摆长之比为 $4:9$，故 A 正确；由于摆长不同，摆动幅度相同，所以摆角不等，故 B 错误；由图可知 $t_{b}$ 时刻，甲、乙两单摆的偏离平衡位置的差值最大，故势能差最大，故 C 正确；$t_{c}$ 时刻甲、乙两单摆都处于平衡位置，由动能定理得 $mgl(1-\cos\theta)=\frac{1}{2}mv^{2}$，设 $A$ 为振幅，$l$ 为摆长，球的位置相对悬挂点的高度差为 $x$，由几何关系知 $\sin\theta=\frac{A}{l}$，联立解得 $v=\sqrt{2gl\left(1-\frac{x}{l}\right)}$，设甲单摆的绳长为 $4l$，则乙单摆的绳长为 $9l$，则 $v_{\text{甲}}=\sqrt{2g\cdot4l\left(1-\frac{x}{4l}\right)}=\sqrt{8gl-2gx}$，$v_{\text{乙}}=\sqrt{2g\cdot9l\left(1-\frac{x}{9l}\right)}=\sqrt{18gl-2gx}$，甲、乙速率不相等，故 D 错误。
}

\item
%% number: 16
%% paperName: 2016年4月浙江选考第 16 题 2 分
%% typeId: 2
%% type: 多选题
%% chapter: 光学
%% point: 光电效应
%% method: 无
%% score: 2
%% degree: 550
%% duplicateId: 0
%% body:
在光电效应实验中，采用极限频率为 $\nu_{c}=5.5\times10^{14}\UHz$ 钠阴极，已知普朗克常量 $h=6.6\times10^{-34}\UJs$，电子质量 $m=9.1\times10^{-31}\Ukg$。用频率 $\nu=7.5\times10^{14}\UHz$ 的紫光照射钠阴极产生光电子的\xzanswer{AD}

\fourchoices[answer=AD]
{动能的数量级为 $10^{-19}\UJ$}
{速率的数量级为 $10^{8}\Ums$}
{动量的数量级为 $10^{-27}\Ukgms$}
{德布罗意波长的数量级为 $10^{-9}\Um$}

%% answer: AD
%% memo:
\memoanswer{
	由光电效应方程得 $E_{\text{km}}=h\nu-W_{0}=h\nu-h\nu_{0}=6.6\times10^{-34}\times(7.5\times10^{14}-5.5\times10^{14})=1.32\times10^{-19}\UJ$，故 A 正确；由 $E_{\text{km}}=\frac{1}{2}mv_{m}^{2}$ 代入数据解得 $v_{m}=5.4\times10^{5}\Ums$，故 B 错误；动量 $P=mv_{m}=9.1\times10^{-31}\Ukg\times5.4\times10^{5}\Ums=4.9\times10^{-25}\Ukgms$，故 C 错误；德布罗意波长 $\lambda=\frac{h}{p}=\frac{6.6\times10^{-34}}{4.9\times10^{-25}}\Um=1.3\times10^{-9}\Um$，故 D 正确。
}

\item
%% number: 17
%% paperName: 2016年4月浙江选考第 17 题 5 分
%% typeId: 4
%% type: 实验题
%% chapter: 功和能
%% point: DIS研究机械能守恒定律
%% method: 无
%% score: 5
%% degree: 650
%% duplicateId: 0
%% body:
\begin{enumerate}
	\item
	在下列学生实验中，需要用到打点计时器的实验有 \tkanswer{BC}（填字母）。

	（A）“探究求合力的方法”　　（B）“探究加速度与力、质量的关系”

	（C）“探究做功与物体速度变化的关系”　　（D）“探究作用力与反作用力的关系”
	\item
	做“验证机械能守恒定律”的实验，已有铁架台、铁夹、电源、纸带、打点计时器，还必须选取的器材是图中的 \tkanswer{BD}（填字母）。
	\item
	某同学在实验过程中，①在重锤的正下方地面铺海绵；②调整打点计时器的两个限位孔连线为竖直；③重复多次实验。以上操作可减小实验误差的是 \tkanswer{②③}（填序号）。
\end{enumerate}

\onepicture[align=right, w=8.5cm]{figs/17.jpg}

%% answer: (1) BC；(2) BD；(3) ②③
\jdanswer{
\begin{enumerate}
	\item
	BC

	\item
	BD

	\item
	②③

\end{enumerate}
}
%% memo:
\memoanswer{
	（1）探究合力的方法，用两根弹簧秤同时拉橡皮筋和一根弹簧秤拉橡皮筋效果相同，探究合力和分力的关系，不需要用到打点计时器，故 A 错误；探究加速度与力、质量的关系，实验中需要通过纸带计算加速度的大小，所以需要打点计时器，故 B 正确；探究做功与物体速度变化的关系，实验中需要通过纸带计算速度，所以需要打点计时器，故 C 正确；探究作用力与反作用力的关系，不需要打点计时器，故 D 错误。

	（2）验证机械能守恒，即验证重物重力势能的减小量和动能的增加量是否相等，所以需要重物，点迹间的时间间隔可以通过打点计时器直接得出，不需要秒表，实验中需要刻度尺测量点迹间的距离，从而得出下降的高度以及瞬时速度的大小，故 BD 正确。

	（3）实验中为了减小实验的误差，调整打点计时器的两个限位孔连线为竖直，从而减小摩擦阻力的影响，重复多次实验，可以减小实验的误差，故应选②③两个操作。
}

\item
%% number: 18
%% paperName: 2016年4月浙江选考第 18 题 5 分
%% typeId: 4
%% type: 实验题
%% chapter: 电路
%% point: 伏安法测电阻
%% method: 无
%% score: 5
%% degree: 550
%% duplicateId: 0
%% body:
在“测绘小灯泡的伏安特性曲线”实验中，小灯泡额定电压为 $2.5\UV$、电流为 $0.3\UA$。
\begin{enumerate}
	\item
	部分连线的实物照片如图 \ref{20164月浙江18a} 所示，请在答题纸上完成实物连接图；
	\item
	某次实验中，当电流表的示数为 $0.18\UA$，电压表的指针如图 \ref{20164月浙江18b} 所示，则电压为 \tkanswer{1.48$\sim$1.50}$\UV$，此时小灯泡的功率是 \tkanswer{0.27}$\UW$；
	\item
	正确测量获得的伏安特性曲线是下列图象中的 \tkanswer{C}（填字母）。
\end{enumerate}

\twopicture[labelA=20164月浙江18a, labelB=20164月浙江18b, align=right, wA=5.0cm, wB=4.5cm]
{figs/18a.png}
{figs/18b.png}

\fourpicture[numA=A, numB=B, numC=C, numD=D, labelprefix={}, align=right, w=2.6cm]
{figs/18c1.png}
{figs/18c2.png}
{figs/18c3.png}
{figs/18c4.png}

%% answer: (1) 实物连接如图所示；(2) 1.48$\sim$1.50，0.27；(3) C
\jdanswer{
\begin{enumerate}
	\item
	实物连接如图所示

	\item
	1.48$\sim$1.50，0.27

	\item
	C

\end{enumerate}
}
%% memo:
\memoanswer{
	（1）小灯泡额定电压为 $2.5\UV$，故电压表选择 $0\sim3\UV$ 的量程；测绘小灯泡的伏安特性曲线时，电压、电流需从 $0$ 开始调节，故滑动变阻器采用分压式接法；电流表采用外接法，故实物连接图如下图所示。

	\onepicture[w=5.0cm]{figs/18_an.jpg}

	（2）结合量程可读出电压表的示数为 $1.50\UV$，故小灯泡的功率为 $P=UI=0.27\UW$。

	（3）随着电压的增大，小灯泡的电阻越来越大，故 C 正确。
}

\item
%% number: 19
%% paperName: 2016年4月浙江选考第 19 题 9 分
%% typeId: 6
%% type: 计算题
%% chapter: 直线运动
%% point: 多段匀变速直线运动
%% method: 无
%% score: 9
%% degree: 650
%% duplicateId: 0
%% body:
如图是上海中心大厦，小明乘坐大厦快速电梯，从底层到达第 119 层观光平台仅用时 $55\Us$。若电梯先以加速度 $a_{1}$ 做匀加速运动，达到最大速度 $18\Ums$，然后以最大速度匀速运动，最后以加速度 $a_{2}$ 做匀减速运动恰好到达观光平台。假定观光平台高度为 $549\Um$。
\begin{enumerate}
	\item
	若电梯经过 $20\Us$ 匀加速达到最大速度，求加速度 $a_{1}$ 及上升高度 $h$；
	\item
	在（1）问中的匀加速上升过程中，若小明的质量为 $60\Ukg$，求小明对电梯地板的压力；
	\item
	求电梯匀速运动的时间。
\end{enumerate}

\onepicture[align=right, h=6.0cm]{figs/19.jpg}

%% answer: (1) $a_{1}=0.9\Umsq$，$h=180\Um$；(2) $F_{N}=654\UN$，方向竖直向下；(3) $t_{0}=6\Us$
\jdanswer{
\begin{enumerate}
	\item
	$a_{1}=0.9\Umsq$，$h=180\Um$

	\item
	$F_{N}=654\UN$，方向竖直向下

	\item
	$t_{0}=6\Us$

\end{enumerate}
}
%% memo:
\memoanswer{
	（1）根据题意，由运动学公式得 $a_{1}=\frac{v_{m}}{t_{1}}=\frac{18}{20}\Umsq=0.9\Umsq$，$h=\frac{1}{2}a_{1}t_{1}^{2}=\frac{1}{2}\times0.9\times20^{2}\Um=180\Um$。

	（2）对小明，由牛顿第二定律得 $F_{N}-mg=ma_{1}$，解得 $F_{N}=654\UN$，由牛顿第三定律知，小明对地板的压力为 $F_{N}'=F_{N}=654\UN$，方向竖直向下。

	（3）设匀速运动时间为 $t_{0}$，运动的总时间为 $t$，作出 $v-t$ 图象，如图所示。由图象可得 $H=\frac{1}{2}(t+t_{0})\times v_{m}$，代入数据解得 $t_{0}=6\Us$。

	\onepicture[w=6.0cm]{figs/19_an.jpg}
}

\item
%% number: 20
%% paperName: 2016年4月浙江选考第 20 题 12 分
%% typeId: 6
%% type: 计算题
%% chapter: 曲线运动
%% point: 竖直平面圆周运动
%% method: 无
%% score: 12
%% degree: 350
%% duplicateId: 0
%% body:
如图所示，装置由一理想弹簧发射器及两个轨道组成。其中轨道 Ⅰ 由光滑轨道 AB 与粗糙直轨道 BC 平滑连接，高度差分别是 $h_{1}=0.20\Um$、$h_{2}=0.10\Um$，BC 水平距离 $L=1.00\Um$。轨道 Ⅱ 由 AE、螺旋圆形 EFG 和 GB 三段光滑轨道平滑连接而成，且 A 点与 F 点等高。当弹簧压缩量为 $d$ 时，恰能使质量 $m=0.05\Ukg$ 的滑块沿轨道 Ⅰ 上升到 B 点；当弹簧压缩量为 $2d$ 时，恰能使滑块沿轨道 Ⅰ 上升到 C 点。（已知弹簧弹性势能与压缩量的平方成正比）
\begin{enumerate}
	\item
	当弹簧压缩量为 $d$ 时，求弹簧的弹性势能及滑块离开弹簧瞬间的速度大小；
	\item
	求滑块与轨道 BC 间的动摩擦因数；
	\item
	当弹簧压缩量为 $d$ 时，若沿轨道 Ⅱ 运动，滑块能否上升到 B 点？请通过计算说明理由。
\end{enumerate}

\onepicture[align=right, w=9.0cm]{figs/20.png}

%% answer: (1) $E_{\text{弹}}=0.1\UJ$，$v_{0}=2\Ums$；(2) $\mu=0.5$；(3) $R_{m}=0.4\Um$，当 $R>R_{m}=0.4\Um$ 时不能到达 B 点
\jdanswer{
\begin{enumerate}
	\item
	$E_{\text{弹}}=0.1\UJ$，$v_{0}=2\Ums$

	\item
	$\mu=0.5$

	\item
	$R_{m}=0.4\Um$，当 $R>R_{m}=0.4\Um$ 时，滑块会脱离螺旋轨道，不能上升到 B 点

\end{enumerate}
}
%% memo:
\memoanswer{
	（1）滑块恰好运动到 B 点时，由机械能守恒得 $E_{\text{弹}}=\Delta E_{k}=\Delta E_{p}=mgh_{1}$，解得 $E_{\text{弹}}=0.1\UJ$。由动能定理得 $\Delta E_{k}=\frac{1}{2}mv_{0}^{2}$，解得 $v_{0}=2\Ums$。

	（2）当弹簧压缩量为 $2d$ 时，由题可知，弹簧的弹性势能是弹簧压缩量为 $d$ 时弹性势能的 $4$ 倍，即 $E_{p2}=4E_{p1}=0.4\UJ$，对滑块从弹簧释放后运动到 C 点的过程，由能量守恒得 $E_{p2}=mg(h_{1}+h_{2})+\mu mg\cos\alpha\cdot L_{BC}$，又 $\mu mg\cos\alpha\cdot L_{BC}=\mu mgL$，联立解得 $\mu=0.5$。

	（3）恰能通过螺旋轨道最高点时，重力提供向心力，须满足的条件是 $mg=m\frac{v^{2}}{R_{m}}$，由机械能守恒得 $\frac{1}{2}mv^{2}=\frac{1}{2}mv_{0}^{2}$，解得 $v=v_{0}=2\Ums$，联立解得 $R_{m}=0.4\Um$。当 $R\leqslant R_{m}=0.4\Um$ 时，滑块能够上升到 B 点；当 $R>R_{m}=0.4\Um$ 时，滑块会脱离螺旋轨道，不能上升到 B 点。
}

\item
%% number: 21
%% paperName: 2016年4月浙江选考第 21 题 4 分
%% typeId: 4
%% type: 实验题
%% chapter: 光学
%% point: 双缝干涉测量波长
%% method: 无
%% score: 4
%% degree: 650
%% duplicateId: 0
%% body:
在“用双缝干涉测量光的波长”实验中，选用红色滤光片和间距为 $0.20\Umm$ 的双缝，双缝与屏的距离为 $600\Umm$。某同学正确操作后，在目镜中看到如图 \ref{20164月浙江21a} 所示的干涉条纹。换成紫色滤光片正确操作后，使测量头分划板刻线与第 $k$ 级暗条纹中心对齐，在目镜中观测到的是图 \ref{20164月浙江21b} 中的 \tkanswer{D}（填字母），此时测量头的读数为 $25.70\Umm$。沿同一方向继续移动测量头使分划板刻线与第 $k+5$ 级暗条纹中心对齐，此时测量头标尺如图 \ref{20164月浙江21c} 所示，其读数是 \tkanswer{19.40}$\Umm$，紫光的波长等于 \tkanswer{420}$\Unm$。

\threepicture[labelA=20164月浙江21a, labelB=20164月浙江21b, labelC=20164月浙江21c, align=right, wA=3.0cm, wB=6.0cm, wC=5.0cm]
{figs/21a.png}
{figs/21b.png}
{figs/21c.png}

%% answer: D，19.40，420
\jdanswer{
D，19.40，420
}
%% memo:
\memoanswer{
	换成紫色滤光片后，光的波长变小，条纹间距变小，故 B 错误，又条纹方向不变，故 C 错误，分划板中心刻线的位置是不变的，A 错误，D 正确。

	游标卡尺的读数为 $19\Umm+8\times0.05\Umm=19.40\Umm$。

	由题意知，相邻条纹的间距为 $\Delta x=\frac{1}{5}\times(25.70-19.40)\Umm=1.26\Umm$，根据 $\Delta x=\frac{L}{d}\lambda$ 解得 $\lambda=4.2\times10^{-7}\Um=420\Unm$。
}

\item
%% number: 22
%% paperName: 2016年4月浙江选考第 22 题 10 分
%% typeId: 6
%% type: 计算题
%% chapter: 磁场
%% point: 带电粒子在磁场中的运动
%% method: 无
%% score: 10
%% degree: 450
%% duplicateId: 0
%% body:
如图为离子探测装置示意图。区域 Ⅰ、区域 Ⅱ 长均为 $L=0.10\Um$，高均为 $H=0.06\Um$。区域 Ⅰ 可加方向竖直向下、电场强度为 $E$ 的匀强电场；区域 Ⅱ 可加方向垂直纸面向里、磁感应强度为 $B$ 的匀强磁场，区域 Ⅱ 的右端紧贴着可探测带电粒子位置的竖直屏。质子束沿两板正中间以速度 $v=1.0\times10^{5}\Ums$ 水平射入，质子荷质比近似为 $\frac{q}{m}=1.0\times10^{8}\UCkg$。（忽略边界效应，不计重力）
\begin{enumerate}
	\item
	当区域 Ⅱ 加电场、区域 Ⅱ 不加磁场时，求能在屏上探测到质子束的外加电场的最大值 $E_{\max}$；
	\item
	当区域 Ⅰ 不加电场、区域 Ⅱ 加磁场时，求能在屏上探测到质子束的外加磁场的最大值 $B_{\max}$；
	\item
	当区域 Ⅰ 加电场 $E$ 小于（1）中的 $E_{\max}$，质子束进入区域 Ⅱ 和离开区域 Ⅱ 的位置等高，求区域 Ⅱ 中的磁场 $B$ 与区域 Ⅰ 中的电场 $E$ 之间的关系式。
\end{enumerate}

\onepicture[align=right, w=7.0cm]{figs/22.png}

%% answer: (1) $E_{\max}=200\UVm$；(2) $B_{\max}\approx5.5\times10^{-3}\UT$；(3) $B=\frac{2E}{v}$
\jdanswer{
\begin{enumerate}
	\item
	$E_{\max}=200\UVm$

	\item
	$B_{\max}\approx5.5\times10^{-3}\UT$

	\item
	$B=\frac{2E}{v}$

\end{enumerate}
}
%% memo:
\memoanswer{
	（1）画出离子的运动轨迹，如图所示。分析轨迹知离子的速度偏转角为 $\theta$，由几何关系得 $\tan\theta=\frac{\frac{H}{2}}{\frac{3L}{2}}$，离子竖直方向的分速度为 $v_{y}=at$，离子通过电场的加速度大小为 $a=\frac{qE_{\max}}{m}$，离子通过区域 Ⅰ 的时间为 $t=\frac{L}{v}$，联立解得 $E_{\max}=200\UVm$。

	\onepicture[w=6.0cm]{figs/22_an1.jpg}

	（2）画出离子的运动轨迹，如图所示。由几何关系知，离子轨迹圆半径满足 $L^{2}+\left(R-\frac{H}{2}\right)^{2}=R^{2}$，解得 $R=\frac{109}{600}\Um$，离子在磁场中运动时，由洛伦兹力提供向心力得 $qvB_{\max}=m\frac{v^{2}}{R}$，解得 $B_{\max}\approx5.5\times10^{-3}\UT$。

	\onepicture[w=5.0cm]{figs/22_an2.jpg}

	（3）画出离子的运动轨迹，如图所示。离子通过区域 Ⅰ 时，速度的偏转角为 $\theta$，由几何关系得 $\tan\theta=\frac{v_{y}}{v}$，由运动学公式得 $v_{y}=at$，$a=\frac{qE}{m}$，$t=\frac{L}{v}$，$v'=\sqrt{v^{2}+v_{y}^{2}}$，离子通过区域 Ⅱ 时，轨迹圆圆心角为 $2\theta$，由几何关系得 $R'\sin\theta=\frac{L}{2}$，由洛伦兹力提供向心力得 $qv'B=m\frac{v'^{2}}{R}$，联立解得 $B=\frac{2E}{v}$。

	\onepicture[w=6.5cm]{figs/22_an3.jpg}
}

\item
%% number: 23
%% paperName: 2016年4月浙江选考第 23 题 10 分
%% typeId: 6
%% type: 计算题
%% chapter: 电磁感应
%% point: 电磁感应电路分析
%% method: 无
%% score: 10
%% degree: 350
%% duplicateId: 0
%% body:
某同学设计了一个电磁推动加喷气推动的火箭发射装置，如图所示。竖直固定在绝缘底座上的两根长直光滑导轨，间距为 $L$。导轨间加有垂直导轨平面向里的匀强磁场 $B$。绝缘火箭支撑在导轨间，总质量为 $m$，其中燃料质量为 $m'$，燃料室中的金属棒 EF 电阻为 $R$，并通过电刷与电阻可忽略的导轨良好接触。

引燃火箭下方的推进剂，迅速推动刚性金属棒 CD（电阻可忽略且和导轨接触良好）向上运动，当回路 CEFDC 面积减少量达到最大值 $\Delta S$，用时 $\Delta t$，此过程激励出强电流，产生电磁推力加速火箭。在 $\Delta t$ 时间内，电阻 $R$ 产生的焦耳热使燃料燃烧形成高温高压气体，当燃烧室下方的可控喷气孔打开后，喷出燃气进一步加速火箭。
\begin{enumerate}
	\item
	求回路在 $\Delta t$ 时间内感应电动势的平均值及通过金属棒 EF 的电荷量，并判断金属棒 EF 中的感应电流方向；
	\item
	经 $\Delta t$ 时间火箭恰好脱离导轨，求火箭脱离时的速度 $v_{0}$；（不计空气阻力）
	\item
	火箭脱离导轨时，喷气孔打开，在极短的时间内喷射出质量为 $m'$ 的燃气，喷出的燃气相对喷气前火箭的速度为 $u$，求喷气后火箭增加的速度 $\Delta v$。（提示：可选喷气前的火箭为参考系）
\end{enumerate}

\onepicture[align=right, h=8.0cm]{\includesvg{figs/23}}

%% answer: (1) $\overline{E}=\frac{B\Delta S}{\Delta t}$，$q=\frac{B\Delta S}{R}$，电流方向向右；(2) $v_{0}=\frac{B^{2}L\Delta S}{mR}-g\Delta t$；(3) $\Delta v=\frac{m'}{m-m'}u$
\jdanswer{
\begin{enumerate}
	\item
	$\overline{E}=\frac{B\Delta S}{\Delta t}$，$q=\frac{B\Delta S}{R}$，电流方向向右

	\item
	$v_{0}=\frac{B^{2}L\Delta S}{mR}-g\Delta t$

	\item
	$\Delta v=\frac{m'}{m-m'}u$

\end{enumerate}
}
%% memo:
\memoanswer{
	（1）由法拉第电磁感应定律得 $\overline{E}=\frac{\Delta\phi}{\Delta t}=\frac{B\Delta S}{\Delta t}$，由电荷量的定义知 $q=\overline{I}\Delta t=\frac{B\Delta S}{R}$，由楞次定律可知，电流方向为 $E\to F$。

	（2）回路中的平均感应电流为 $\overline{I}=\frac{\overline{E}}{R}=\frac{B\Delta S}{R\Delta t}$，棒 EF 受到的平均安培力为 $\overline{F}=B\overline{I}L$，设竖直向上为正，由动量定理得 $(\overline{F}-mg)\Delta t=mv_{0}$，联立解得 $v_{0}=\frac{B^{2}L\Delta S}{mR}-g\Delta t$。

	（3）以火箭为参考系，设竖直向上为正，由动量守恒得 $-m'u+(m-m')\Delta v=0$，解得 $\Delta v=\frac{m'}{m-m'}u$。
}

\end{enumerate}

\end{document}
